\section{X Files and Bigoted Opinions}

Due to the flare up of a chronic condition, I realized how difficult it is to think while in physical discomfort. Not that I missed it much ... So I put that down time to good use by studying several manifestations of popular culture. This post is the best I can do with such a poor starting point.

\begin{quotationx}
I read a depressingly inane magazine article by a Logical Positivist ... The burden of his teaching seems to be this:

``Since we cannot really say anything about anything, let us be content to talk about the way in which we say nothing.'' That is an excellent way to organize futility.

After all, even nothingness has its dignity: but here not even the dignity of nothingness is respected. There must be the mechanical clicking of the thought machine manufacturing nothing about nothing, as if even nothing had at all costs to be organized, and presented as if it were something. As if it had to be talked about.

The atheist existentialist has my respect: he accepts his honest despair with stoic dignity. And despair gives his thought a genuine content, because it expresses an experience --- his confrontation with emptiness. But those others confront only the mechanical output of their own thinking machine. They don't have the imagination or the good sense to stand in awe at real emptiness. In fact, their rationalizations seem to be a complacent evasion: as if logical formulas somehow could give something to stand on in the abyss.

And now: just wait until they start philosophizing with computers!
\flright{\textsc{Thomas Merton}, \emph{Conjectures of a Guilty Bystander}}
\end{quotationx}


\paragraph{X Files Philosophy}


In the TV Show, \emph{X-Files}, the two agents hold radically different metaphysical positions. The female character, Dana Scully, is a committed logical positivist, certain that there is no reality beyond what material science can prove. Fox Mulder, on the other hand, believes in paranormal phenomena, conspiracy theories, and so on.

Although Mulder's point of view is not a valid alternative to positivism, we could let it stand as strictly entertainment. However, a somewhat serious thinker was recently brought to our intention. He claims to be an idealist, but by that he means little more than a more systematic and sophisticated version of Mulder's philosophy. Specifically, he seems to treat philosophy as a scientific hypothesis except, in his case, he regards consciousness as independent of matter and hence a causative factor in the material world. Hence, he accepts all sorts of ``psychic powers'' as ``proofs'' of his philosophy.

Unfortunately, that is not at all what we mean by ``philosophical idealism''. Rather, the fundamental premise is that the world can be known by thinking about it; hence, it must be rational. This is the opposite of scientific positivism, which sees the world as the result of random forces, so that even thinking itself arises from non-rational causes.

With 7 billion people on the planet, why is there not one of them who can guess the 4 digit code to unlock that iPhone via psychic powers?

\paragraph{What he said}


In \emph{Conjectures of a Guilty Bystander}, \textbf{Thomas Merton} tells the story of a Governor's visit to the Monastery. The Guv said, ``You monks know you cannot be happy because you have material possessions.'' Certainly, it was poorly phrased, so Merton decided to take a literal reading of it. A novice monk protested, ``But that is not what he \emph{meant}.'' Merton drew this lesson:

\begin{quotex}
How true it is that everyone instinctively pays attention not to what a politician actually says, but to what he seems to want to say.
\end{quotex}


Gornahoor draws a different lesson:

\begin{quotex}
How true it is that everyone instinctively pays attention not to what a politician actually says, but to what they would like him to say.
\end{quotex}


\paragraph{Bigoted Opinions}


I'm amazed at the emotional involvement that some people have for their preferred candidate, even though there is little or no benefit for them personally. But then again, democracy is the worst form of government, except for those that will follow afterward (to paraphrase someone). That's bad news for those hoping for a reversal of current trends somehow.

That's because in a democracy, everyone is entitled to his own opinion. But if those opinions are simply random products of evolution\footnote{\url{https://www.gornahoor.net/?p=8474}}, what value have they? Obviously, such opinions may have practical reality if they work to hold society together. Yet that is unsatisfactory to most people since they are convinced that their opinions are factually real. I'm afraid you can't eat your Darwinism and have it, too, which what the educated class seems to want.

The real answer is nihilism. The consequence of nihilism is that there are only bigoted opinions vying with each other. I'll leave you a longish text by \textbf{Harry Neumann}\footnote{\url{https://www.gornahoor.net/?page_id=8496}} for you to ponder without comment. Read it as Merton reads the atheist existentialist and develop contempt for those complacent evaders.

\paragraph{The New Religion}


I discovered a channel on satellite radio that euphemistically goes by the name the ``Catholic Channel''. However, to a large degree, it is unlike any religion I know. For example, there was a discussion about whether burying a statue of St Joseph in your yard would help you sell your house.

I've noticed that there has developed a cottage industry of converts and reverts who tend to dominate blogging. Unfortunately, their enthusiasm cannot overcome their ignorance of Tradition. They lack that sense of Tradition that lingers, even in apostates like Santayana and many others. They tend to jumble things together, treating all teachings as having equal importance. I have to assume that RCIA classes are quite unhelpful.

One such woman today was talking about gluttony and Oprah Winfrey's battle with fat. Of course, Oprah has a ``food addiction'' since ``gluttony must be a very rare sin''. In a more sane era, the hostess would have been yanked off the air. What exactly is a food addiction, medically? It describes a behavior, not a condition. Sin is hard to talk about today, but it entails an inner bondage to a behavior. In other words, a food addiction. More on this later, but for now I would recommend \emph{Sin Revisited} by \textbf{Solange Hertz}.

\paragraph{Filling the Pews}


On another show a Jesuit was waxing ecstatic about how great the Church would become with woman deacons who could give sermons on issues of concern to them. For good measure, he suggested that acceptance of homosexuality would also be nice. Now with estimates of up to 50\% of priests of that orientation, I can't imagine it is a problem.

I've been told that most parishes are already dominated by women, so what is the issue? The real problem is the lack of young men in the pew, for obvious reasons. That Jesuit talked about encounter and other squishy words and used ``privilege''\footnote{\url{https://en.wikipedia.org/wiki/Privilege\_(social\_inequality)}} as a verb. He can't imagine how off-putting his language is for most men.

\paragraph{Ted Cruz and Tradition}


Ted Cruz belongs to a religious movement called the New Apostolic Reformation\footnote{\url{https://www.youtube.com/results?search\_query=New+Apostolic+Reformation}}. They are very Old Testament oriented with menorahs, stars of David, and shofars used in their services. As a sort of ape of Tradition, they seek to restore the offices of Priest and King, except they have a very democratic view of those roles. Their ultimate goal is dominion of the culture in its seven domains: religion, family, education, government, media, arts \& entertainment, and business. Ultimately, they will manifest preternatural powers on the earth.

In the abstract, that sounds traditional. They seem serious, with a few million adherents, and a plan. Compare that to NRx\footnote{\url{http://neoreaction.net/}}, which consists of a few guys meeting in a bar. Logic is just one of the liberal arts and it is ineffective without rhetoric. Logic is actually easier, so much more time should be spent on rhetorical skills. The first step is to learn to recognize it.

\paragraph{The Sense of Wonder}


In esoterism, there is the distinction between knowledge and understanding. On the one hand, there is the mere accumulation of facts and information, the other is grasping the whole which links the facts together. Idealist philosophy is ``knowledge of the whole'', not the development of super powers.

Ultimately, it is a matter of ``seeing''. You can't argue with the complacent evader, but you may get him to ``see''. If you don't see, then spend some quiet time and think about existence. Boris Mouravieff writes:

\begin{quotex}
Exterior man feels the transcendent character of the Love arising from the Absolute I and Absolute II. The first in essence reaches his consciousness in the form of a perception of existence. But he considers the beauty of the Universe and its life as mere information, instead of as a marvelous gift capable of constantly stimulating his sense of wonder and kindling his gratitude.
\end{quotex}

\textit{update 18 Feb 2021}

I guess I got the worldview of Dana Scully incorrect. I should have spent more of my life watching TV.

I forget who is that philosopher of psychic powers.

Neo-reaction is totally dead now and the links don't work. Does anyone know what happened?

Plato was the ``someone'' who predicted that tyranny follows democracy. So don't be surprised.

I don't know who is the fellow in the turban or what he would think of the comparison.


\flrightit{Posted on 2016-02-18 by Cologero}

\begin{center}* * *\end{center}

\begin{footnotesize}
\begin{sffamily}
\texttt{Harun on 2016-02-19 at 05:25 said:}

Despite Nietzsche disgust for christian tradition is philosophy holds deep metaphysical truth, even Suzuki roshi mentioned him as the european closer to a zen understanding.

Most people who think they are nihilist are just radical relativists who still evade the question. The positivist replace his faith in God with faith in the scientific method, when he realize how even that method can't produce truth he turns into a relativists. From that point nothing in the world has any meaning, the relativist is a sophist who holds only the opinion that suits him at any moment.

Many so called nihilist refuse any moral code to live as they want, they think they are closer to the ubermensch but in fact they are what Nietzsche called ``the last man''.

The last man is actually the result of negative nihilism, is the product of the ``age of absolute sinfulness'' as described by the idealist Fichte.

Positive nihilism is something closer to a Nigredo: abandoning all ideas, beliefs, attachments and operate a painful clean to ourselves.

It's a step toward initiatic death, I feel pity for modern ``nihilist'' who think Nietzsche just told them to live like animals and satisfy their lower needs without any moral concern.

Unfortunately it's only normal that the profane see the opposite of truth: too many times gnostic doctrines were interpretated not to ascend but to descen into the lower realms of animal istincts.

Guénon explains too well how an initiation can be perverted and become it's opposite.

\hfill

\texttt{Dart on 2016-02-19 at 20:12 said:}

Scully is actually a Roman Catholic.

\hfill

\texttt{James on 2016-02-21 at 08:42 said:}

\url{https://40.media.tumblr.com/14fbaea6ff4021399e1c61539ead6207/tumblr\_nj329srMKC1tx9o6mo1\_1280.jpg}


Hi, Cologero, I just thought that this man resembles you so much.

\hfill

\texttt{Cologero on 2016-02-21 at 22:33 said:}

I'm flattered, James.

\hfill

\texttt{Dennis on 2016-02-22 at 17:07 said:}

While Scully is skeptical of Mulder's more outlandish claims regarding aliens and conspiracy theories, it's not clear that she believes there is ``no reality beyond what material science can prove.'' She's seen too much in her years with Mulder to take such a dogmatic logical positivist view. She is a lapsed Catholic in whom some residual faith seems to linger -- she wears a Cross around her neck, for example -- seems open to the possibility of the divine (and of aliens!), and she has never been identified as an atheist.

\hfill

\texttt{Logres on 2016-02-23 at 01:58 said:}

That's an interesting point about Logic \& Rhetoric. Rhetoric seems to carry a logic of its own, almost three or four dimensional. Some of the Southern writers like Mel Bradford were trying to say something quite similar, when they looked at the founding tradition, Lincoln's speeches, etc. Perhaps to someone with a clear enough head and clean enough heart, Rhetoric offers its own course in Wisdom and higher Logic, because ``language is sermonic''.

\hfill

\texttt{Aidan on 2021-02-18 at 07:31 said:}

Neoreaction slowly petered out even though Curtis Yarvin is more active now and Urbit is picking up steam. I was in personal contact with a few that led to a surreal project.

While the likeness isn't quite uncanny it's there, he would probably be quietly amused.

\hfill

\texttt{Greg on 2021-02-21 at 23:27 said:}

Only the Orthosphere seems active in reaction. Social matters I think drove the movement and when that site went dead so did neoreaction.

\hfill

\texttt{Balrog Booger on 2022-02-18 at 12:38 said:}

Neoreaction went the way so many things on the internet do: as soon as the term started to pick up steam it became \emph{cringe,} and everybody just stopped calling themselves that. The ``movement'' as such still exists, but has migrated from written blogs to the easier to digest mass medium of YouTube. Rather than orbiting around Curtis Yarvin, it increasingly is moving toward such video essayists as The Distributist, yet another guy on the internet calling himself Charlemagne, Auron MacIntyre, and their general sphere of social contacts. These days the term of choice for such people seems to be ``dissident rightist.'' Calling yourself a dissident seems insane to me, but the wind of fashion blows in whatever direction it will.

\end{sffamily}
\end{footnotesize}

